\documentclass[ocean-paper.tex]{subfiles}
\begin{document}
\section{Human Games}\label{appendix:humans}
%\subsection{Matches against skilled human players}
% Worth talkng abut MMR?
% \dota uses a matchmaking rating (MMR) system in order that awards points to players each game based on a combination of the win or loss outcome and how favored a team was in the matchup. As \openaifive progressed, we competed against increasingly stronger teams made up of higher MMR players. At the highest level, professional teams tend to not be ranked by individual MMR but rather historical success as a team.

\begin{table}
\begin{tabular}{c|c|c|c|l}
Opponent & Result & Duration & Version & Restrictions\\
\hline
\multicolumn{5}{l}{June 6, 2018 - Internal Event}\\
\hline
Internal team & win & 15:15 (surr) & 7.13 & Mirror match, multiple couriers, no invis\\
Internal team & win & 20:51 & 7.13 & Mirror match, multiple couriers, no invis\\
Audience team & win & 31:33 & 7.13 & Mirror match, multiple couriers, no invis\\
Audience team & win & 23:33 (surr) & 7.13 & Mirror match, multiple couriers, no invis\\
\hline
\multicolumn{5}{l}{August 5, 2018 - Benchmark}\\
\hline
Caster team & win & 21:38 (surr) & 7.16 & Drafted, multiple couriers\\
Caster team & win & 24:56 (surr) & 7.16 & Drafted, multiple couriers\\
Caster team & lose & 35:47 & 7.16 & Audience draft, multiple couriers\\
\hline
\multicolumn{5}{l}{August 9, 2018 - Private eval}\\
\hline
Team Secret & win & 17:00 (surr) & 7.16 & Drafted, multiple couriers\\
Team Secret & lose & 48:46 & 7.16 & Drafted, multiple couriers\\
Team Secret & lose & 38:55 & 7.16 & Drafted, multiple couriers\\
\hline
\multicolumn{5}{l}{August 22-23, 2018 - The International}\\
\hline
Pain Gaming & lose & 52:29 & 7.19 & Pre-set lineup\\
Chinese Legends & lose & 45:44 & 7.19 & Pre-set lineup\\
\hline
\multicolumn{5}{l}{October 5, 2018 - Private eval}\\
\hline
Team Lithium & win & 48:57 & 7.19 & TI pre-set lineup\\
Team Lithium & win & 48:16 & 7.19 & TI pre-set lineup\\
Team Lithium & win & 31:33 & 7.19 & Drafted\\
\hline
\multicolumn{5}{l}{January 16, 2019 - Private eval}\\
\hline
SG Esports & win & 24:29 (surr) & 7.19 & TI pre-set lineup\\
SG Esports & win & 25:08 (surr) & 7.19 & Drafted\\
SG Esports & win & 27:36 (surr) & 7.20 & Mirror match\\
SG Esports & win & 25:30 (surr) & 7.20 & Mirror match\\
\hline
\multicolumn{5}{l}{February 1, 2019 - Private eval}\\
\hline
Alliance & win & 17:11 & 7.20d & Drafted\\
Alliance & win & 31:33 & 7.20d & Drafted\\
Alliance & win & 28:16 & 7.20d & Reverse drafted\\
\hline
\multicolumn{5}{l}{April 13, 2019 - \openaifinals}\\
\hline
OG & win & 38:18 & 7.21d & Drafted\\
OG & win & 20:51 & 7.21d & Drafted\\
\end{tabular}
\caption{Major matches of \openaifive against high-skill human players.}
\label{table:milestonegames}
\end{table}

See \autoref{table:milestonegames} for a listing of the games \openaifive played against high-profile teams.

%\subsection{Arena: benchmarking against the community}
% \label{appendix:arena}
% In addition to a few short games against the very best players, we sought to determine if the \openaifive could be consistently exploited by creative or out-of-distribution play. On April 18-21 2019, we launched a system called {\it \openaifive Arena}, in which thousands players from around the world competed against \openaifive simultaneously. 3,193 teams competed in 7,257 games (in addition to 35,466 ``cooperative games,'' where less than 5 human players allied with \openaifive agents on the same team). 
% Of these \openaifive won\footnote{Human players often abandoned losing games rather than playing them to the bitter end, even abandoning games right after an unfavorable hero selection draft before the main game begins. 
% \openaifive does not abandon games, so we count abandoned games as wins for \openaifive. These abandoned games (3140 of the 7215 wins) likely includes a small number of games that were abandoned for technical or personal reasons.} 7215, or 99.4\%. 

% To determine whether teams could reliably beat \openaifive by finding and exploiting weaknesses, we tracked whether teams could win multiple games in a row. 
% \todo{Jonathan: the next part seems a bit too wordy, worth trimming}
% For example, if OpenAI beat a team ten times, then lost once, then won the next five, we do not consider that team to have found an exploitable strategy. 
% However if OpenAI won 10 games to a team and then lost 2 games, then the team stopped playing, we consider that that team likely found a repeatable exploit at the end.
% Of the 3,193 teams which participated, 22 were able to win at all, and only 4 were able to win multiple games in a row. 
% All 4 of those later lost a game, indicating that perhaps they had not found a foolproof exploit. 
% One team won 10 games in a row, by exploiting a weakness in \openaifive's understanding of invisible heroes. 
% This team then lost their next game, demonstrating that even taking advantage of the exploit requires careful execution and is far from simple.
\end{document}